\chapter{Application and Case Study}
\label{chap:case_study}

This chapter gives a clear overview of the dataset and experimental
process used in this study. It describes the evaluation dataset with
examples, explains how the 786 model responses were collected, reports on
data quality checks that found real bugs, and shows where to find the
full code and repository for reproducibility.

\section{Data Introduction}
\label{sec:data-intro}

The evaluation dataset has 131 math questions from two sources: Hendrycks
MATH (40 Easy and 40 Medium questions) and AIME (51 Hard questions). Each
question includes a unique ID, category, difficulty tier, source, the
question text, and the correct final answer.
Table~\ref{tab:examples-full} gives a simplified preview of the
processed dataset: one example question and answer from each of the four
mathematical categories (Algebra, Combinatorics, Number Theory,
Probability), covering one Easy, one Medium, and two Hard questions, to
show what a typical question and its answer look like across the
dataset.

\begin{longtable}{p{1.6cm}p{1.5cm}p{6.7cm}p{1.5cm}}
  \caption{Simplified preview of the processed dataset: one example question and answer from each category, covering one Easy, one Medium, and two Hard questions.} \label{tab:examples-full} \\
  \toprule
  Question ID & Difficulty & Question & Answer \\
  \midrule
  \endfirsthead
  \toprule
  Question ID & Difficulty & Question & Answer \\
  \midrule
  \endhead
  \bottomrule
  \endfoot
  ALG\_E\_003 & Easy &
    $3^n = 3 \cdot 9^3 \cdot 81^2$. What is the value of $n$? & 15 \\[4pt]
  COMB\_M\_002 & Medium &
    A diagonal of a polygon is a segment joining two nonconsecutive
    vertices of the polygon. How many diagonals does a regular octagon
    have? & 20 \\[4pt]
  NT\_H\_008 & Hard &
    Find the sum of all positive integers $n$ such that $n+2$ divides
    the product $3(n+3)(n^2+9)$. & 49 \\[4pt]
  PROB\_H\_001 & Hard &
    Jen enters a lottery by selecting 4 distinct elements of $S =
    \{1,2,\ldots,10\}$. Then four elements of $S$ are drawn at random.
    Jen wins a prize if at least two of her numbers were drawn, and
    wins the grand prize if all four of her numbers were drawn. The
    probability that Jen wins the grand prize given that Jen wins a
    prize is $m/n$ where $m$ and $n$ are relatively prime positive
    integers. Find $m+n$. & 116 \\
\end{longtable}

These four examples show the difference between the Hendrycks-sourced
Easy/Medium tiers and the AIME-sourced Hard tier. The Algebra (Easy)
example is a single equation to solve for one unknown. The
Combinatorics (Medium) example is a direct counting formula. The Number
Theory (Hard) example requires reasoning about a divisibility condition
across every positive integer, not a fixed set of candidates, and the
Probability (Hard) example requires a probability calculation over a
combinatorial system with several derived quantities before the final
step. The two Hard-tier examples show the kind of open-ended,
multi-step derivation characteristic of AIME, compared with the more
directly computable Hendrycks-sourced examples at the Easy and Medium
tiers.

\section{Experimental Pipeline}
\label{sec:pipeline}

The experimental pipeline gathers 786 responses in total. This comes from
running 131 questions with both Chain-of-Thought and Tree-of-Thought
methods using three models (131 $\times$ 2 $\times$ 3 = 786). Instead of
processing all in one long, uninterruptible batch, the pipeline is built
to be resumable and uses checkpoints. This was important because
free-tier API quotas were reached several times during data collection.

Each experiment is identified by its Question ID, model, and prompt type.
Before starting, the pipeline checks which experiment combinations are
already in the results file and skips those that are done. This means
the same script can simply be re-run to proceed from wherever it
previously stopped, without manual observation or restarting from
scratch. Each result is saved to the results file as soon as it
finishes, one row at a time. If the run is interrupted by an error,
quota limit, or if the process is stopped, all completed work is
preserved.

This approach was important because the three providers have very
different free-tier rate limits. Mistral is limited to 2 requests per
minute, so it is much slower than Gemini or Groq. For example, Mistral
takes about 5 to 6 hours for a full dataset pass, while the other two
take only tens of minutes. To avoid wasting time, all three models run
at the same time, each in its own thread. This means the total run time
depends only on the slowest model. If a model's daily quota runs out
during a run, its thread stops without affecting the others or losing
any collected data. The run can be repeated later when the quota
resets, and only the unfinished combinations are processed.

Each result row also includes provenance information in addition to the
raw response. This includes the exact model version returned by the
provider's API (which may differ from the configured model alias, as
explained in Chapter~\ref{chap:methodology}), the particular prompt
version used, and the start time, end time, and execution time for each
call.

\section{Data Quality Assurance}
\label{sec:data-quality}

To ensure the automated pipeline works correctly, a clear rule is
enforced: a response's Answer Correct field and its Error Type must
always agree with each other. If a response is marked incorrect but
labelled with the error type ``Correct'' (or the other way around), it
signals a bug in the pipeline, not a real reasoning issue. Rather than
assuming this rule holds, the code checks every batch of results
against it before analysis begins, and stops immediately if a mismatch
is found.\footnote{Implemented in the function
\texttt{assert\_no\_stale\_error\_labels}, in \texttt{src/utils/io.py}.}
This check surfaced two real answer-extraction bugs during error
classification that would have otherwise gone unnoticed and produced
wrong data.

The first bug was related to LaTeX fraction notation. The extraction
logic could handle the basic \texttt{\textbackslash frac\{a\}\{b\}}
macro, but not the \texttt{\textbackslash dfrac} (display-style) or
\texttt{\textbackslash tfrac} (text-style) versions. For example, in
question PROB\_E\_003 (Groq, ToT), the model's response used
\texttt{\textbackslash dfrac\{3\}\{16\}} as its final answer; the macro
stayed in the text, and the fallback numeric extraction picked up just
the digit ``3'' instead of the correct value ``3/16.'' This meant a
correct answer was marked as incorrect purely because of a parsing
error, and it was fixed by updating the logic to recognise
\texttt{\textbackslash dfrac} and \texttt{\textbackslash tfrac} as well.

The second bug was about embedded ordinal numbers. The extraction logic
searched for numbers in the response, but in a sentence like ``The 158th
marble is gray,'' the number 158 is just an ordinal reference, not the
answer. For example, in question NT\_E\_009 (Groq, CoT), the real
answer, ``gray,'' is a word, but because the logic matched ``158''
before checking for words, the pipeline recorded ``158'' as the answer
instead. This was fixed by ensuring the logic does not treat numbers
with ordinal suffixes (like ``158th'') as standalone numeric answers.

Regression tests were added for both fixes, using the exact failing
inputs described above.\footnote{The regression tests are implemented
in \texttt{tests/test\_response\_parser.py}.} This way, if future
changes to the extraction logic bring back either bug, the tests will
catch them automatically instead of letting them slip by.

\section{Reproducibility}

The full source code for the experimental pipeline, analysis scripts,
and statistical tests described in this dissertation is available at
\url{https://github.com/AnushreeM2701/LLM_Project}. This includes the
answer evaluator and statistical test implementations, the full
verbatim CoT and ToT prompt templates, and the complete dataset
composition and question listing.
